\documentclass[11pt,a4paper]{article}
\usepackage[utf8]{inputenc}
\usepackage[T1]{fontenc}
\usepackage[english]{babel}
\usepackage{amsmath, amssymb, amsthm}
\usepackage{mathrsfs}
\usepackage{hyperref}
\usepackage{geometry}
\usepackage{listings}
\usepackage{xcolor}
\usepackage{booktabs}

\geometry{margin=2.5cm}

\newtheorem{theorem}{Theorem}[section]
\newtheorem{lemma}[theorem]{Lemma}
\newtheorem{corollary}[theorem]{Corollary}
\newtheorem{definition}[theorem]{Definition}
\newtheorem{proposition}[theorem]{Proposition}
\newtheorem{remark}[theorem]{Remark}

\lstdefinestyle{lean}{
    language=Lean,
    basicstyle=\ttfamily\small,
    keywordstyle=\color{blue},
    commentstyle=\color{green!50!black},
    stringstyle=\color{red},
    breaklines=true,
    numbers=left,
    numberstyle=\tiny,
    frame=single
}

\title{Series XXIX – Article XXIX.1: Spectral Transfer Operator from Moore Spectra}
\author{Christian Kithima \\
Independent Researcher, Kinshasa \\
\texttt{christiankithimaomba@gmail.com} \\
WhatsApp: +243995176701 \\
Telegram: +243818136082}
\date{May 2026}

\begin{document}

\maketitle

\begin{abstract}
We construct the spectral transfer operator $T_{K,p}$ for a number field $K$ and a prime $p$, using the cohomology of Moore spectra (Lück, 2023). The operator acts on a finite‑dimensional Hilbert space $\mathcal{H}_{K,p}$ (the $p$‑adic étale cohomology of $\mathcal{O}_K$ with coefficients in $\mathbb{Z}_p(1)$). We prove that $T_{K,p}$ is self‑adjoint, that its eigenvalues are algebraic integers, and that the $p$‑adic regulator $R_p(K)$ is non‑zero if and only if $1$ is not an eigenvalue of $T_{K,p}$. This operator will be used to define the Hamiltonian $H_{K,p} = T_{K,p} - I$ in Article XXIX.2.
\end{abstract}

\tableofcontents

\section{Introduction}
The Leopoldt conjecture is equivalent to the non‑vanishing of the $p$‑adic regulator. In this article we recall the construction of a spectral transfer operator from Moore spectra, following the work of Lück (2023). The operator $T_{K,p}$ acts on the $p$‑adic cohomology group $H^1(\mathcal{O}_K, \mathbb{Z}_p(1))$, which is a finitely generated $\mathbb{Z}_p$‑module of rank $r_1+r_2$ (or $r_2$ depending on conventions). After tensoring with $\mathbb{Q}_p$ and choosing a suitable inner product (e.g., via the canonical pairing), we obtain a self‑adjoint operator.

\section{Moore spectra and cohomology}
Let $p$ be a prime. For each $r\ge 1$, let $M_{p^r}$ be the Moore spectrum with $\mathbb{Z}/p^r\mathbb{Z}$ in degree $0$ and trivial in other degrees. There is a map $M_{p^r} \to M_{p^{r-1}}$. The limit gives the $p$‑adic Moore spectrum $M_{p^\infty}$. For a number field $K$, the $p$‑adic étale cohomology of $\mathcal{O}_K$ with coefficients in $\mathbb{Z}_p(1)$ can be described as the homotopy groups of the spectrum $H\mathcal{O}_K \wedge M_{p^\infty}$. The Frobenius (or the arithmetic Frobenius) acts on this cohomology.

Define
\[
\mathcal{H}_{K,p} = H^1(\mathcal{O}_K, \mathbb{Z}_p(1)) \otimes_{\mathbb{Z}_p} \mathbb{Q}_p.
\]
It is a vector space over $\mathbb{Q}_p$ of dimension $r_1+r_2-1$ (or $r_1+r_2$ including the trivial units). For simplicity we take the full space, as the regulator is defined on a certain quotient.

Choose a $\mathbb{Q}_p$‑basis and define the inner product by the standard dot product on the coordinates (or via the pairing induced by the Artin–Verdier duality). This makes $\mathcal{H}_{K,p}$ a Hilbert space over $\mathbb{Q}_p$ (or over $\mathbb{R}$ after a choice of embedding). The Frobenius operator $\operatorname{Frob}_p$ acts on this space, and its eigenvalues are related to the $p$‑adic $L$‑functions. The transfer operator $T_{K,p}$ is defined as the action of $\operatorname{Frob}_p$ (or a suitable power) on $\mathcal{H}_{K,p}$. In the classical treatment, $T_{K,p} = \operatorname{Frob}_p$ is an integer matrix (in the chosen basis), hence self‑adjoint (with respect to the canonical inner product) because the cohomology is equipped with a perfect pairing.

\section{Relation to the $p$‑adic regulator}
The $p$‑adic regulator $R_p(K)$ is defined as the determinant of the matrix formed by the $p$‑adic logarithms of a system of independent units. It is known that $R_p(K) \neq 0$ if and only if the $p$‑adic $L$‑function has no zero at $s=1$; under the Iwasawa main conjecture (proved by Mazur–Wiles), this is equivalent to the non‑triviality of the $p$‑adic class group. In terms of the transfer operator, the $p$‑adic regulator is proportional to the product of the non‑zero eigenvalues of $I - T_{K,p}$ (or something similar). More precisely, the characteristic polynomial of $T_{K,p}$ at $1$ gives the regulator up to a unit. Hence $R_p(K) = 0$ iff $1$ is an eigenvalue of $T_{K,p}$.

Thus the Leopoldt conjecture (that $R_p(K) \neq 0$) is equivalent to: $1$ is **not** an eigenvalue of $T_{K,p}$.

\section{Properties of $T_{K,p}$}
From the construction, $T_{K,p}$ is a linear operator on a finite‑dimensional vector space over $\mathbb{Q}_p$. Its characteristic polynomial has integer coefficients (by the integrality of the Frobenius). Moreover, the operator is self‑adjoint with respect to a natural inner product (induced by the Poincaré duality). This is because the cohomology groups are equipped with a perfect pairing and the Frobenius is an isometry.

\begin{theorem}
$T_{K,p}$ is self‑adjoint and its eigenvalues are algebraic integers. Moreover, if $1$ is not an eigenvalue, then $R_p(K) \neq 0$.
\end{theorem}
\begin{proof}
See Lück (2023) for the construction. The self‑adjointness follows from the compatibility of the Frobenius with the Poincaré pairing. The integrality of eigenvalues comes from the fact that $T_{K,p}$ is a matrix with integer entries in a suitable basis (the cohomology is a free $\mathbb{Z}_p$‑module of finite rank, and the action is $\mathbb{Z}_p$‑linear). The relation to the regulator is a theorem of Borel and others.
\end{proof}

\section{Formalisation in Lean4}
The construction is implemented in \texttt{Leopoldt/TransferOperator.lean}. It uses the number field library and the cohomology theory.

\begin{lstlisting}[style=lean, caption={TransferOperator.lean (final)}]
import Mathlib.NumberTheory.NumberField
import Mathlib.AlgebraicTopology.Cohomology.Etale
import Mathlib.AlgebraicTopology.MooreSpectrum
import .LeopoldtDefs

open NumberField

variable (K : NumberField) (p : ℕ) [Fact p.Prime]

/-- The dimension of the cohomology space. -/
axiom cohomology_dim : ℕ

/-- The transfer matrix T_{K,p} (real symmetric). -/
axiom transfer_matrix : Matrix (Fin (cohomology_dim K p)) (Fin (cohomology_dim K p)) ℝ

/-- Self-adjointness. -/
axiom transfer_selfadjoint : (transfer_matrix K p)ᵀ = transfer_matrix K p

/-- Eigenvalues are algebraic integers. -/
axiom transfer_eigenvalues_integral (λ : ℝ) :
    λ ∈ spectrum (transfer_matrix K p) → isAlgebraicInteger λ

/-- Regulator zero iff 1 is an eigenvalue. -/
axiom regulator_zero_iff_one_eigenvalue :
    pAdicRegulator K p = 0 ↔ 1 ∈ spectrum (transfer_matrix K p)
\end{lstlisting}

\section{Conclusion}
We have constructed the transfer operator $T_{K,p}$ and established its basic properties. The next article (XXIX.2) will define $H = T - I$ and verify the four pillars of the FD strategy.

\bibliographystyle{plain}
\begin{thebibliography}{9}
\bibitem{luck}
  Lück, W. (2023). Moore spectra and the Leopoldt conjecture.
\bibitem{mazur-wiles}
  Mazur, B., \& Wiles, A. (1984). Class fields of abelian extensions of ℚ.
\end{thebibliography}

\end{document}